%%%%%%%%% begin snippet
%% You need to add the package "tabularx".
%% Place the snippet right after \begin{document}

% need tabularx
\documentclass{article}
\usepackage{tabularx}
\usepackage{graphicx}
\usepackage{pdfpages}
\usepackage{listings}

\begin{document}
    \begin{titlepage}
           \begin{center}
                 \begin{huge}
                       %% Update assignment number here
                       \textbf{Assignment 2}
                 \end{huge}
           \end{center}

           \begin{center}
                 \begin{large}
                       Machine Learning 1, SS23
                 \end{large}
           \end{center}

           \begin{center}
                \begin{tabularx}{\textwidth}{|>{\hsize=.33\hsize}X|>{\hsize=.33\hsize}X|>{\hsize=.33\hsize}X|}

                       \hline
                       \multicolumn{3}{|c|}{\textbf{Team Members}} \\
                       \hline
                       Last name & First name & Matriculation Number \\
                       \hline
                       Grassl & Ifeoma & 12011965 \\
                       \hline
                       Royer & Christoph & 12004184 \\
                       \hline

                \end{tabularx}
           \end{center}

    \end{titlepage}

    \section{Neural Networks}
    \subsection{PCA and Classification}
    \subsubsection{PCA for dimensionality reduction}
    \textit{For implementation see } \texttt{nn\_classification.py}.

    We used $320$ components, which gave a ratio of $95.05\%$ explained variance.

    \subsubsection{Varying the number of hidden neurons}
    \textit{For implementation see } \texttt{nn\_classification.py}.

    \textbf{Output of the program:}
    \begin{lstlisting}
Number of hidden neurons: 2
Train accuracy: 0.5913. Test accuracy: 0.4262
Loss: 1.0635
Number of hidden neurons: 10
Train accuracy: 0.9982. Test accuracy: 0.7215
Loss: 0.0156
Number of hidden neurons: 100
Train accuracy: 1.0000. Test accuracy: 0.7724
Loss: 0.0054
Number of hidden neurons: 200
Train accuracy: 1.0000. Test accuracy: 0.7918
Loss: 0.0044
    \end{lstlisting}

    Underfitting is generally easy to spot, as it manifests in poor performance both on training and on test data.
    This is often the case if the model does not have enough capacity (ie. not enough complexity); in this case, the loss converges on a bad value.
    The model cannot improve past a (relatively bad) optimum.

    Overfitting takes place when a powerful model is trained on too little data.
    In this case, the model learns the training data -- including its inaccuracies -- too well.
    The side effect of this is that the model fails to predict test data accurately because of an overfixation on the particularities of the training dataset.

    We could not detect overfitting in this example.
    Thus we chose the model with $200$ hidden neurons, as it had the best test accuracy.

    \subsubsection{Preventing overfitting}
    Our preliminary guess was that model \textbf{c} -- with \textit{alpha} = $0.1$ and \textit{early\_stopping} = True -- would work best.
    The reasoning was that multiple prevention mechanisms would be best at catchinng overfitting.

    It turns out that model \textbf{a} -- with only \textit{alpha} = $0.1$ -- worked best.

    \pagebreak
    We can see the test output here -- note that Variation 0 (model \textbf{a}) improved on the test accuracy from the previous task:

    \begin{lstlisting}
Number of hidden neurons: 200
Variation 0:
Train accuracy: 1.0000. Test accuracy: 0.8378
Loss: 0.0678
Variation 1:
Train accuracy: 0.9770. Test accuracy: 0.7700
Loss: 0.0541
Variation 2:
Train accuracy: 0.9709. Test accuracy: 0.7772
Loss: 0.1722
    \end{lstlisting}

    \subsubsection{Variability of the performance}
    \textit{For implementation see } \texttt{nn\_classification.py}.

    We chose model \textbf{a}, as it showed the best output.
    Note that this is a form of overfitting, as we are choosing hyperparameters based on their performance on the test set,
    but the assignment requires this.

    The implementation outputs the following:
    \begin{lstlisting}
Seed: 12011965
Train accuracy: 1.0000. Test accuracy: 0.8354
Loss: 0.0668
Seed: 12004184
Train accuracy: 1.0000. Test accuracy: 0.8475
Loss: 0.0656
Seed: 1234
Train accuracy: 1.0000. Test accuracy: 0.8402
Loss: 0.0672
Seed: 7
Train accuracy: 1.0000. Test accuracy: 0.8402
Loss: 0.0673
Seed: 4321
Train accuracy: 1.0000. Test accuracy: 0.8329
Loss: 0.0671
Train accuracy overall:
On the train set: 1.0000 +/- 0.0000 [1.0000:1.0000]
On the test set: 0.8392 +/- 0.0050 [0.8329:0.8475]
    \end{lstlisting}

    Changing the seed changes the pseudo-random initial state of the MLP.
    This should not change the optimal solution,
    but a fortunate choice of the initial state may drastically reduce the number of iterations needed for convergence.

    \subsubsection{Loss curve}
    We chose seed $12004184$, as it showed the best performance in the previous test.
    Note that this is again a form of overfitting, but the assignment requires this.

    \includegraphics[width=\textwidth]{code/plots/losscurve}

    \pagebreak
    \subsubsection{Prediction on the test set}
    The following is the output of the trained model's prediction and the confustion matrix on the test set:

    \begin{lstlisting}
Predicting on the test set
        precision recall    f1-score    support

0       0.77      0.85      0.81        40
1       0.83      0.91      0.87        44
2       0.85      0.73      0.79        48
3       0.79      0.79      0.79        39
4       0.82      0.89      0.86        37
5       0.80      0.84      0.82        43
6       0.79      0.68      0.73        44
7       0.94      0.86      0.90        37
8       0.71      0.80      0.75        30
9       0.98      0.96      0.97        51

accuracy                           0.83       413
macro avg       0.83      0.83      0.83       413
weighted avg       0.84      0.83      0.83       413

[[34  0  0  0  2  2  1  0  0  1]
 [ 1 40  0  1  0  0  2  0  0  0]
 [ 0  2 35  3  1  3  0  0  4  0]
 [ 1  0  0 31  1  0  2  0  4  0]
 [ 0  1  0  0 33  1  1  0  1  0]
 [ 1  2  1  0  1 36  2  0  0  0]
 [ 4  2  2  3  0  3 30  0  0  0]
 [ 3  0  1  0  0  0  0 32  1  0]
 [ 0  0  2  1  2  0  0  1 24  0]
 [ 0  1  0  0  0  0  0  1  0 49]]

    \end{lstlisting}
    \textit{Recall} measures the ratio of true positive in relation to the amount of samples that were available given class.
    It thus measures how many samples were missed to be recognised as members of the class.
    We can easily recalculate this for class $i$ with $\frac{tp_i}{support_i}$, with $tp_i$ being the $i$th diagonal entry in the confusion matrix (= true positive classification).

\end{document}

%%%%%%%%% end snippet
